\documentclass{article}
\usepackage[utf8]{inputenc}
\usepackage{amsmath}
\usepackage{geometry}
\geometry{a4paper, margin=1in}

\title{Mathematical Formulations of Solution Quality KPIs}
\author{Nurse Scheduling System (NSS)}
\date{\today}

\begin{document}

\maketitle

\section{Roster \& Staffing Efficiency}

\subsection{Demand Coverage}
Measures the percentage of total patient demand met by the planned schedule (Regular and Overtime shifts).
\begin{equation}
\text{Coverage} = \left( \frac{\sum_{i \in I} \sum_{j \in J} \sum_{k \in K} (sr_{ijk} + so_{ijk})}{\text{Total Demand}} \right) \times 100\%
\end{equation}

\subsection{Understaffing Rate (Emergency Reliance)}
Measures the reliance on expensive emergency staff (recourse actions) to meet demand across all scenarios.
\begin{equation}
\text{Understaffing Rate} = \left( \frac{\sum_{j \in J} \sum_{k \in K} \sum_{w \in W} p^w \alpha_{jkw}}{\sum_{j \in J} \sum_{k \in K} \sum_{w \in W} p^w (sr_{ijk} + so_{ijk} + \alpha_{jkw})} \right) \times 100\%
\end{equation}

\subsection{Overtime Ratio}
Measures the proportion of overtime usage relative to regular working hours.
\begin{equation}
\text{Overtime Ratio} = \left( \frac{\sum_{i \in I} \sum_{j \in J} \sum_{k \in K} so_{ijk}}{\sum_{i \in I} \sum_{j \in J} \sum_{k \in K} sr_{ijk}} \right) \times 100\%
\end{equation}

\section{Economic Efficiency}

\subsection{Cost Per Shift}
The blended average cost to staff a single shift, incorporating regular, overtime, and emergency wages.
\begin{equation}
\text{Cost Per Shift} = \frac{\text{Objective Value (Total Expected Cost)}}{\sum_{j \in J} \sum_{k \in K} \sum_{w \in W} p^w (\sum_{i \in I}(sr_{ijk} + so_{ijk}) + \alpha_{jkw})}
\end{equation}

\section{Biological Risk (Fatigue)}

\subsection{Fatigue Rate}
Percentage of shifts performed by nurses exceeding the safety threshold ($F_{threshold} = 0.5$).
\begin{equation}
\text{Fatigue Rate} = \left( \frac{|\{ (i,j) : F_{ij} > 0.5 \}|}{\text{Total Assignments}} \right) \times 100\%
\end{equation}

\subsection{Average Workforce Fatigue}
The mean fatigue level across the entire team for the full roster period.
\begin{equation}
\bar{F} = \frac{1}{N \times J} \sum_{i=1}^{N} \sum_{j=1}^{J} F_{ij}
\end{equation}

\subsection{Peak Fatigue}
The maximum fatigue level reached by any individual nurse in the schedule.
\begin{equation}
F_{\text{peak}} = \max_{i \in I, j \in J} \{ F_{ij} \}
\end{equation}

\section{Workforce Equity}

\subsection{Workload Balance (Standard Deviation)}
Measures the fairness of shift distribution among the nurse pool.
\begin{equation}
\sigma = \sqrt{\frac{\sum_{i=1}^{N} (x_i - \mu)^2}{N}}
\end{equation}
Where $x_i$ is the total shifts for nurse $i$, and $\mu$ is the average shifts per nurse.

\end{document}
